\section{Experiment Results}



\begin{table}[]
    \centering
    \caption{Pass@1 comparison between proposed AnalogSAGE and baselines.}
    \label{tab:baseline_vs_full}
    \resizebox{\linewidth}{!}{
        \begin{tabular}{cccccccccccc}
            \hline
            \textbf{Framework} & \textbf{1} & \textbf{2} & \textbf{3} & \textbf{4} & \textbf{5} & \textbf{6} & \textbf{7} & \textbf{8} & \textbf{9} & \textbf{10} & \textbf{Avg.}\\
            \midrule
            \textbf{SPICEPilot \cite{vungarala2024spicepilot}}       & 0.0 & 0.0 & 0.0 & 0.0 & 0.0 & 0.0 & 0.0 & 0.0 & 0.0 & 0.0 & 0.0\\
            \textbf{AnalogCoder-Pro \cite{lai2025analogcoder}}  & 0.0 & 0.0 & 0.0 & 20.0 & 0.0 & 0.0 & 0.0 & 0.0 & 0.0 & 0.0 & 2.0 \\
            \hline
            \textbf{Artisan-style SAGE}          & \textbf{100.0} & 80.0 & 80.0 & 60.0 & \textbf{100.0} & 20.0 & 0.0 & 0.0 & 20.0 & 0.0 & 46.0 \\
            \textbf{AnalogSAGE} & \textbf{100.0} & \textbf{100.0} & \textbf{100.0} & \textbf{100.0} & \textbf{100.0} & \textbf{80.0} & \textbf{100.0} & \textbf{100.0} & \textbf{100.0} & \textbf{80.0} & \textbf{96.0}\\
            \hline
        \end{tabular}}
\end{table}


\begin{table}[]
    \centering
    \caption{Evaluation result under four configurations.}
    \label{tab:ablation_study_splithead}
    \resizebox{\linewidth}{!}{
    \begin{tabular}{cccccccccccc}
        \hline
        \textbf{Configuration} & \textbf{Metric} &
        \textbf{1} & \textbf{2} & \textbf{3} & \textbf{4} & \textbf{5} &
        \textbf{6} & \textbf{7} & \textbf{8} & \textbf{9} & \textbf{10} \\
        \hline

        \multirow{1}{*}{W/o IO and EM} & Pass@1
            & \textbf{100.0} & 80.0 & \textbf{100.0} & 40.0 & 80.0
            & 40.0 & 20.0 & 60.0 & 40.0 & 0.0 \\

        \multirow{1}{*}{W/o IO} & Pass@1
            & \textbf{100.0} & \textbf{100.0} & \textbf{100.0} & 80.0 & \textbf{100.0}
            & 20.0 & 40.0 & 40.0 & 60.0 & 0.0 \\

        \multirow{2}{*}{W/o EM} & Pass@1
            & \textbf{100.0} & \textbf{100.0} & \textbf{100.0} & \textbf{100.0} & \textbf{100.0}
            & 40.0 & 60.0 & 80.0 & \textbf{100.0} & 20.0 \\

        & Avg Iter
            & \textbf{1.0} & 1.2 & 1.4 & 1.4 & \textbf{1.2}
            & NA & NA & NA & 1.8 & NA \\

        \multirow{2}{*}{Full(AnalogSAGE)} & Pass@1
            & \textbf{100.0} & \textbf{100.0} & \textbf{100.0} & \textbf{100.0} & \textbf{100.0}
            & \textbf{80.0} & \textbf{100.0} & \textbf{100.0} & \textbf{100.0} & \textbf{80.0} \\

        & Avg Iter
            & \textbf{1.0} & \textbf{1.0} &\textbf{1.2} & \textbf{1.0} & \textbf{1.2}
            & NA & \textbf{1.2} & \textbf{1.4} & \textbf{1.0} & NA \\
        \hline
    \end{tabular}}
\end{table}

\subsection{Performance comparison}

Table~\ref{tab:baseline_vs_full} compares the proposed AnalogSAGE against baselines. SPICEPilot, which relies solely on zero-shot prompting without simulation feedback or iterative reasoning, fails on all tasks, yielding a 0\% pass rate and 0\% Pass@1. AnalogCoder-Pro introduces limited feedback but lacks grounded numerical optimization under the PDK. It occasionally discovers a correct design after many trials, achieving a 10\% pass rate and only a 2\% Pass@1. Artisann-style SAGE benefits from predefined circuit templates and simulation-guided parameter tuning, achieving a 70\% pass rate. However, its first-try success remains below 50\%, demonstrating weak generalization to strict specifications. In contrast, AnalogSAGE achieves a \textbf{100\% pass rate}, meaning every design task is eventually solved, and a \textbf{96.0\% Pass@1}, meaning 96\% of tasks succeed without requiring additional trials. This reflects a \textbf{10 $\times$} pass rate and a \textbf{48$\times$} in Pass@1 over the prior works, demonstrating stronger reliability, faster convergence, and better generalization across difficulty levels.

\subsection{Role of self-evolving memory layers}


To quantify the contribution of the proposed self-evolving memory design, we conducted experiments under four configurations.  
The lowest configuration disables both the \emph{Introspective Optimization Layer (IO)} and the \emph{Evolution Memory Layer (EM)}, denoted as \emph{W/o IO and EM}.  
The second removes only IO (\emph{W/o IO}), while the third removes only EM (\emph{W/o EM}).  
The full system keeps both layers enabled. All configurations were evaluated across ten design tasks, and the results are summarized in Table~\ref{tab:ablation_study_splithead}.


The system without either IO or EM can handle easy and medium tasks, but completely fails on the most challenging ones. It also exhibits unstable convergence, often requiring multiple trials to obtain a feasible solution.
Adding only EM improves performance on medium tasks by enabling cross-task knowledge transfer, which initializes the search closer to feasible operating regions. The gradual Pass@1 improvement (from 20\% to 60\% in Tasks~6--9) confirms that accumulated design priors benefit later tasks, although the system still fails on the hardest case due to the lack of adaptive correction.
In contrast, enabling only IO allows the system to eventually solve Task~10, demonstrating that introspective self-correction can repair failing design trajectories. However, its Pass@1 remains low (20\%), indicating that without insight priors the system lacks stability.

When both IO and EM are enabled, the framework achieves the best overall performance, reaching a 100\% pass rate and the highest Pass@1 with fewer iterations. This shows that the \emph{Evolution Memory Layer} provides informed initialization for new design tasks, while the \emph{Introspective Optimization Layer} adaptively corrects reasoning paths using grounded feedback. The two layers are complementary and jointly necessary for robust autonomous analog design.


\subsection{Parameter optimization efficiency}



% Beyond reasoning accuracy, we further examined the efficiency of AnalogSAGE on parameter optimization performance, which would predict parameter ranges for each transistor and passive component by the LLM that directly determine the efficiency of the subsequent numerical optimization as a more accurate and tighter parameter range leads to a smaller search space and faster convergence, and a higher likelihood of success.
Beyond reasoning accuracy, we also evaluated AnalogSAGE’s efficiency in parameter optimization. In this stage, the LLM predicts parameter ranges for each transistor and passive component. These predicted ranges directly shape the search space for the subsequent numerical optimizer. Tighter and more accurate ranges reduce the search space, accelerate convergence, and increase the likelihood of success.
Table~\ref{tab:sizing_LLM} shows the normalized parameter search space across different configurations where the smaller is better.  
Our AnalogSAGE achieves an average space of 0.26, indicating a \textbf{4$\times$ reduction} in parameter search space, which outperforms other three configurations.  
These results confirm that the \emph{Knowledge Module} provides physically consistent parameter bounds, while the \emph{Introspective Optimization Layer} further tightens feasible regions through reflective, simulation-driven feedback.
\begin{table}[]
    \centering
    \caption{Comparison of normalized parameter search space across different configurations.}
    \label{tab:sizing_LLM}
\resizebox{\linewidth}{!}{
\begin{tabular}{cccccccccccc}
    \hline
    \textbf{Method} &
    \textbf{1} & \textbf{2} & \textbf{3} & \textbf{4} & \textbf{5} &
    \textbf{6} & \textbf{7} & \textbf{8} & \textbf{9} & \textbf{10} & \textbf{Avg.} \\
    \hline

    W/o IO and Knowledge
        & 0.11 & 0.36 & 0.67 & 1.00 & 1.00
        & 1.00 & 0.62 & 0.31 & 0.37 & 0.29 & 0.57  \\

    W/o IO
        & 0.29 & 0.03 & 0.89 & \textbf{0.31} & \textbf{0.37}
        & 0.35 & 0.36 & 0.58 & 0.06 & \textbf{0.0008} & 0.32 \\

    W/o Knowledge
        & \textbf{0.03} & 0.37 & 0.64 & 1.00 & 1.00
        & 1.00 & 1.00 & \textbf{0.30} & 0.37 & 0.33 & 0.60 \\
    Full (AnalogSAGE)
        & 0.07 & \textbf{0.02} & \textbf{0.27} & 1.00 & 0.45
        & \textbf{0.32} & \textbf{0.05} & 0.37 & \textbf{0.04} & 0.003 & \textbf{0.26} \\

    \hline
\end{tabular}}
\end{table}
